\section*{Resumen}

La administración de una tienda especializada en libros, discos y artículos de colección requiere mantener información coherente sobre el catálogo, las existencias, las ventas y la recepción de mercancía. La gestión distribuida en registros independientes dificulta consultar la disponibilidad y reconstruir las operaciones realizadas. Para atender esta necesidad se desarrolló \textbf{Archivum: Bibliotheca \& Discos}, una aplicación web que integra un catálogo público con herramientas administrativas de inventario y ventas.

La aplicación está implementada con Laravel 13 sobre PHP 8.4 y organiza su información en una base de datos MariaDB 11.2. Docker Compose coordina los servicios de aplicación PHP-FPM, base de datos y servidor web Nginx. Las vistas Blade y los recursos de interfaz se compilan con Vite y Tailwind CSS. El modelo de datos relaciona productos y sus especializaciones para libros, discos y coleccionables con variantes SKU, existencias por sucursal, piezas únicas, movimientos, carritos, reservas, pedidos, pagos, envíos y devoluciones.

Entre las funciones disponibles se encuentran la búsqueda y el filtrado del catálogo, el registro y autenticación de usuarios, la recuperación de contraseña, la administración de proveedores y productos, la recepción de mercancía, los ajustes y traspasos de inventario, el seguimiento de pedidos, la gestión de devoluciones y la exportación de reportes. El flujo de compra web es simulado y aplica reservas temporales de cinco minutos. El catálogo curado se puede sincronizar mediante un comando Artisan y un seeder versionado. La ejecución de la suite automatizada concluyó con 190 pruebas aprobadas y 1\,335 afirmaciones; no se realizaron mediciones de carga ni pruebas concurrentes.

\vspace{0.4cm}
\noindent\textbf{Palabras clave:} gestión de inventario, comercio electrónico, Laravel, MariaDB, Docker, catálogo especializado.
